\section{Communication}


\begin{frame}{Project Governance Report}
    \begin{itemize}
        \item The `constitution' of a project
        \item A contract between stakeholders
        \item The actual plan, typically comprising
        \begin{itemize}
            \item user requirements ('need' or reason)
            \item detailed solution design
            \item stakeholder analysis
            \item work breakdown structure 
            \item definition of work packages and responsibilities
            \item project plan: network / BPMN plan
            \item project plan: temporal plan / Gantt diagram
            \item governance plan: corporation, controls
            \item identity and roles
            \item communication strategy 
            \item documentation, controls and reporting
            \item risk register
        \end{itemize}
    \end{itemize}
\end{frame}



\newcommand{\mybox}[1]{\fbox{\parbox[t]{2.5cm}{#1}}}
\begin{frame}{Perspectives on the Vision}  % depends on: forest

    
    \showmytikz{../shared/tikz/tikz-PM-vision-perspectives}
    
 
    \vfill
    Source: following \cite[p.113]{Roberts}
\end{frame}

\begin{frame}{Communication Strategy}
    \begin{tabular}{llllll}
        \toprule
         Stakeholder & Deliverable & Frequency & Owner & Type & $\dots$ \\
         \midrule
         Sponsors       & report & monthly & PM & presentation & $\dots$ \\
         Tech team 1    & meeting & weekly & PM & coordination & $\dots$ \\
         Sales team     & update & bi-weekly & PM & key facts & $\dots$ \\
         Users    & mailout & monthly & sales & briefing & $\dots$ \\
         $\dots$ & $\dots$ & $\dots$ & $\dots$ & $\dots$ &  $\dots$ \\
         \bottomrule
    \end{tabular}
\end{frame}

\begin{frame}{Communication Strategy}        
    \begin{itemize}
        \item Strategy per phase differently
        \item Information `points' into Gantt-diagram
        \item Schedule communication
        \item Communication is often \textbf{a milestone}
        \item Communication \textbf{flows along organisational chart} (?)
        \item \textbf{Communication is key!}
        \item Communciation is \textbf{not a one-way task!}
    \end{itemize}
    
\end{frame}




% \newcommand{\mybox}[1]{\fbox{\parbox{2.5cm}{\centering #1}}}
\begin{frame}{Communication in Project Management}
    \showmytikz{../shared/tikz/tikz-PM-communication-directions}
    \par
    Following \cite[p.~205]{KerznerHarold2017Pm:a}
\end{frame}


\renewcommand{\mybox}[1]{\fbox{
    \begin{minipage}[t][10ex][t]{2cm}
        \centering #1
    \end{minipage}
}}
% \begin{frame}{Customer Communciations}
%     \begin{tikzpicture}[scale=0.7]\small
%         \def\myx{3}
%         \def\myy{3.5}
%         \node       at (-\myx,\myy+5.5cm) {CUSTOMER};
%         \node       at (\myx,\myy+5.5cm) {CONTRACTOR};
%         \node (sp1) at (-\myx,\myy) {\mybox{Project Sponsor or Executive}};
%         \node (pm1) at (-\myx,0) {\mybox{Project Manager or Contracts Officer}};
%         \node (em1) at (-\myx,-\myy) {\mybox{Employees}};
        
%         \node (sp2) at (\myx,\myy) {\mybox{Project Sponsor or Executive}};
%         \node (pm2) at (\myx,0) {\mybox{Project Manager or Contracts Officer}};
%         \node (em2) at (\myx,-\myy) {\mybox{Employees}};
        
%         \draw[<->] (sp1) -- (pm1);
%         \draw[<->] (pm1) -- (em1);

%         \draw[<->] (sp2) -- (pm2);
%         \draw[<->] (pm2) -- (em2);
        
%         \draw[<->,dashed] (sp1) -- node[above] {informal} (sp2) ;
%         \draw[<->] (pm1) -- node[above] {\textbf{formal}} (pm2) ;
%         \draw[<->,dashed] (em1) -- node[above] {informal} node[below] {\emph{trouble?}} (em2) ;
%     \end{tikzpicture}
%     \cite[p.205]{KerznerHarold2017Pm:a}
% \end{frame}






% \begin{frame}{Information transmission (Shannon/Weaver)}

%     The most simple view is technial:
    
%     \medskip

%     \begin{tikzpicture}
%         \node (S) at (0,0) {\fbox{Sender}};
%         \node (I) at (3,0) {\fbox{Information}};
%         \node (R) at (6,0) {\fbox{Receiver}};
%         \draw[->] (S) -- (I);
%         \draw[->] (I) -- (R);
%     \end{tikzpicture}
    
        
%     \begin{itemize}
%         \item `Magic Bullet', `Hypodermic needle'
%         \item `Information Deficit Model':
%         \item Assumes receiver `needs to know'
%       %  \item \emph{But: }  who is active? The sender, or the receiver, or both? $\rightarrow$ Uses and gratification theory

%     \end{itemize}   
%     Source: \cite{ShannonClaudeElwood1964Tmto} etc.
    
%     \framebreak
    
%     \begin{itemize}
%         \item Is this \emph{information} or merely \emph{data}?
%         \item Information can be measured, and this measure is the same as for \emph{entropy}
%         \item Information cannot be \textbf{compressed}
%         \item Information is \textbf{data with meaning}
%         \item But what is \textbf{meaning?}
%         \item When is communication \textbf{meaningful?}
%     \end{itemize}

%     \framebreak
    
    
%     \begin{tikzpicture}
%         \usetikzlibrary{arrows};
%         \node (S) at (0,0) {\fbox{Sender}};
%         \node (T) at (2,0) {\fbox{Transmitter}};
%         \node (s1) at (4,-1) {distortion};
%         \node (C) at (4,0) {\fbox{Channel}};
%         \node (R) at (6,0) {\fbox{Receiver}};
%         \node (D) at (8,0) {\fbox{Destination}};
%         \draw[->] (S) -- (T);
%         \draw[->] (T) -- (C);
%         \draw[->] (C) -- (R);
%         \draw[->] (R) -- (D);
%         \draw[->] (s1) -- (C);
%         \draw[->,bend left=45] (D) to node[above] {feedback} (S);
%     \end{tikzpicture} 
    
%     Source: e.g. \cite{ShannonClaudeElwood1964Tmto} 
% \end{frame}



\begin{frame}{Information Economics - Principal-Agent-Problem }
    \showmytikz{../shared/tikz/tikz-PM-principal-agent}
    \end{frame}

\begin{frame}{Information Economics - Principal-Agent-Problem }    \begin{minipage}{0.7\textwidth}
    
        \begin{itemize}
            \item The principal needs some external good or service which is beyond her own expertise.
        
            \item The agent delivers this good or service, but has far more expertise about the solution, its making, and so forth, than the principal. 
        
            \item There is an \textbf{information bias.} The result is \textbf{inefficiency.}
            
            \item This results in higher prices and lower volumes.
        
        \end{itemize}
    
    \end{minipage}
    \par\medskip
    Source (example for energy efficiency): \cite{OgdenCharlesKay1985Tmom}

\end{frame}

\begin{frame}{Adverse selection, moral hazard}

    Some examples: 
    
    \medskip
    
    \begin{minipage}{0.7\textwidth}
        \begin{enumerate} 
            \item the home-builder knows less than the architect less than the actual site-manager...
            \item risk-prone seek insurance, risk-adverse don't
            \item bad products need advertisement
        \end{enumerate}
    \end{minipage}
    \par\medskip
    Source (classic): \cite{GeorgeA.Akerlof1970TMfL}
\end{frame}

    
% \begin{frame}{Principal-Agent-Problem in the Electricity Sector}

%     \fbox{\parbox{0.4\textwidth}{
%     Principal: \textbf{The Public}
%     \begin{itemize}
%         \item `I want electricity cheap'
%         \item `I want electricity green'
%         \item `I want electriciy reliably'
%     \end{itemize}
%     }}
%     $\leftrightarrow$
%     \fbox{\parbox{0.4\textwidth}{
%     Agent: \textbf{Energy provider (?)}
%     \begin{itemize}
%         \item `Sector coupling...'
%         \item `Smart meters...'
%         \item `There is a trade-off between...'
%     \end{itemize}
%     }}
    
% \end{frame}



% \subsection{Knowledge Gaps II: Oneself}


% \begin{frame}{Difficulties to self-estimate}
%     \begin{itemize}
%         \item Information asymmetries exist between people
%         \item These are difficult to identify
%         \item Identification requires research (!)
%         \item Dunning-Kruger effect: inability to judge how complex / complicated something unknown is leads to excessive self-confidence
%         \begin{itemize}
%             \item `I drive better than the average driver'
%             \item `I am not a good guitar player'
%             \item `everyone can paint something like that'
%             \item `why is nobody interested in my project?'
%         \end{itemize}
%     \end{itemize}
%     \vfill
%     Literature tipp: \cite{kahneman2011thinking}
% \end{frame}

\begin{frame}{Four Stages of Competence}
    \showmytikz{../shared/tikz/tikz-PM-four-stages-competence}
    \vfill
    Source: own design / commons
\end{frame}

\begin{frame}{'Rumsfeld'-Diagramm}
    \showmytikz{../shared/tikz/tikz-PM-rumsfeld-matrix}
    \par\medskip
    [own design / commons]
\end{frame}

\begin{frame}{Planning is Information Processing}
    \begin{itemize}
        \item You plan something \textbf{novel}
        \item There are necessarily \textbf{unknowns}
        \item You must start with \textbf{research}
        \item Research is \textbf{asking questions}
        \item You involve \textbf{people}
        \item You must \textbf{communicate}
        \item Knowledge is a \textbf{resource}
        \item People \textbf{must know what they must konw}
        \item \textbf{How do you know what someone knows?}
    \end{itemize}
\end{frame}

% \begin{frame}{Systems communicate}
%     \begin{tikzpicture}
%         \draw (0,0) rectangle (2,2) ;
%         \draw (3,2) rectangle (5,0) ;
%         \draw[->] (2,0.5) -- (3.2,0.5);
%         \draw[<-] (2,1.5) -- (3,1.5);
%         \draw (3.2,0.2) rectangle (4,1);
%     \end{tikzpicture}
%     \par
%     \cite{ParsonsTalcott1951Tss}, \cite{luhmann1984soziale}, \cite{RopohlGunter1979ESdT}
    
%     \framebreak
    
%     \begin{tikzpicture}[scale=0.5]
%         \path[fill=blue] (0,-2) -- (0,-3) -- (0.75,-3) -- (0.75,-5) -- (5.25,-5) -- (5.25,-4) -- (6,-4) -- (6,-2) -- (4,-2) -- (4,-4) -- (4.75,-4) -- (4.75,-4.5) -- (1.25,-4.5) -- (1.25,-3) -- (2,-3) -- (2,-2) -- cycle;
%         \path[draw,thick] (0,-1) -- (0,-3) -- (0.75,-3) -- (0.75,-5) -- (5.25,-5) -- (5.25,-4) -- (6,-4) -- (6,-1);
%         \path[draw,thick] (4,-1) -- (4,-4) -- (4.75,-4) -- (4.75,-4.5) -- (1.25,-4.5) -- (1.25,-3) -- (2,-3) -- (2,-1);
%     \end{tikzpicture}

% \end{frame}


% \begin{frame}{Planning Communication}

%     \begin{tikzpicture}
%       \draw (-1,-1.5) rectangle (7,4.5);
%       \node[draw, diamond,fill=red] (m1) at (1,1.5) {};
%       \draw (1.5,0) rectangle (3,3);
%       \node[draw, diamond,fill=red] (m2) at (3.5,1.5) {};
%       \draw (4,0) rectangle (5.5,3);
%       \node[draw, diamond,fill=red] (m3) at (6,1.5) {};
%       \draw[->] (m3) -- (5.5,1.5);
%       \draw[->] (4,1.5) -- (m2);
%       \draw[->] (m2) -- (3,1.5);
%       \draw[->] (1.5,1.5) -- (m1);
%       \node (def) at (0,4) {definition};
%       \node (cont) at (3.5,-1) {control};
%       \draw[->] (def) -| (m2);
%       \draw[->] (cont) -- (m2);
%       \draw[->, dashed] (4,3) -- (def);
%       \draw[->] (def) -- (1.5,3);
%       \draw[->] (cont) -- (3,0);
%     \end{tikzpicture}
    
%     Definition comes before the milestone,
    
%     control after the milestone.
    
%     \framebreak
    
%     \begin{itemize}
%         \item You must define milestones
%         \item You must communicate milestones
%         \item Milestones are communicated
%     \end{itemize}
    
%     \framebreak
    
%     %\begin{adjustbox}{width=\textwidth}
%         \smartdiagramset{back arrow disabled=true}
%         \smartdiagram[flow diagram:vertical]{understand, formulate, communicate}
%     %\end{adjustbox}
    
%     \framebreak
    
%     PDCA Cycle (Plan - do - check - ajust)
    
%     \smartdiagram[circular diagram:clockwise]{Plan,Do,Check,Adjust}
    
%     \framebreak
    
%     \smartdiagram[flow diagram:vertical]{discuss, agree, communicate}
    
% \end{frame}

% \begin{frame}{Making a Decision}
%     \begin{itemize}
%         \item Your job is to \textbf{identify} the best solution
%         \item and to \textbf{commit everyone} to it
%         \item Modern manager is a \textbf{problem solver}
%         \item Most problems are information / communication based
%         \item Formulate s.m.a.r.t.\ goals
%     \end{itemize}
% \end{frame}



\begin{frame}{Planning Communication}
    \begin{itemize}
        \item Processes: physical work flows: \textbf{what}
        \item Financial, material flows: \textbf{with what}
        \item People: \textbf{who}
        \item People must \textbf{know and love their roles}
        \item \textbf{Communication must be planned}
        \item Give room for communication
    \end{itemize}
\end{frame}



\begin{frame}{Information in Gantt Diagram}
      \begin{ganttchart}[y unit chart=0.4cm]{1}{9}
      \gantttitlelist{1,...,9}{1} \\
      \ganttgroup{Project cabin}{1}{9} \\
      \ganttbar[bar/.append style={fill=red}]{Plan the cabin (4)}{1}{4} \\
      \ganttbar{Site management}{2}{9} \\
      \ganttgroup{Carpenters}{2}{5} \\
      \ganttbar{Procure Wood (1)}{2}{2} \\
      \ganttbar{Construct (3)}{3}{5} \\
      \ganttgroup{Utility Guys}{3}{7} \\
      \ganttbar{Plan Utilities (1)}{3}{4} \\
      \ganttbar{Procure Utilities (1)}{5}{5} \\
      \ganttbar{Utilities (2) }{6}{7}\\
      \ganttgroup{Interior Guys}{4}{9} \\
      \ganttbar{Plan Equip (1)}{4}{4}\\
      \ganttbar{Procure Equip (1)}{5}{7}\\
      \ganttbar{Equip (2)}{8}{9}
      \ganttvrule{}{4}
    \end{ganttchart}
    \end{frame}

\begin{frame}{Information in Gantt Diagram}    
    \begin{ganttchart}[y unit chart=0.4cm]{1}{12}
      \gantttitlelist{1,...,12}{1} \\
      \ganttgroup{Project cabin planning}{1}{7} \\
        \ganttbar{First design}{2}{2} \\
        \ganttbar{Discuss with architect(s)}{3}{4} \\
        \ganttbar{Discuss with contractors}{5}{6} \\
        \ganttbar{Discuss with authorities}{4}{6} \\
        \ganttbar{Discuss with users}{1}{3} 
        \ganttbar{Discuss with users}{5}{8} \\
        \ganttgroup{Execution}{8}{12} \\
        \ganttbar{Start building}{7}{9} \\
        \ganttbar{Start equipping}{9}{9}
        \ganttlink{elem1}{elem2}
        \ganttlink{elem2}{elem3}
        \ganttlink{elem3}{elem8}
        \ganttlink{elem6}{elem9}
    \end{ganttchart}
    \end{frame}

\begin{frame}{Information in Gantt Diagram}    
    \begin{itemize}
        \item Communication \textbf{is key}
        \item Communication is \textbf{involvement}
        \item \textbf{Planning is communication}
    \end{itemize}
    
\end{frame}

% \begin{frame}{Scope Creep}
%     \begin{itemize}
%         \item Sponsor realises he wants more
%         \item Management realises product needs more
%         \item Milking the sponsor or milking the mananger?
%         \item What is necessary?
%         \item How can this all be paid for?
%     \end{itemize}
% \end{frame}

% \begin{frame}{Product Description}
%     \begin{itemize}
%         \item Title
%         \item Purpose
%         \item Composition
%         \item Derivation
%         \item Format
%         \item Audience
%         \item Quality criteria
%         \item Quality method
%     \end{itemize}
% \end{frame}

% \begin{frame}{The Three Constraints}
%     \begin{tikzpicture}[scale=0.8]
%         \coordinate (A) at (0,0);
%         \coordinate (B) at (-3,-5);
%         \coordinate (C) at (3,-5);
%         \coordinate (X) at (0,-3.35);
%         \draw (A) -- (B) node[midway,above,rotate=60] {Costs};
%         \draw (B) -- (C) node[midway,below] {Quality/Performance/Technology};
%         \draw (C) -- (A) node[midway,above,rotate=-60] {Time};
%         \path[draw=blue, -stealth, rotate=90,
%       postaction={decorate,decoration={text effects along path,
%         text={\ Within good customer relations \ }, text align=center, 
%         text effects/.cd, 
%           text along path, 
%           every character/.style={fill=white, yshift=-0.5ex}}}]
%       (X) circle (4);
%     \end{tikzpicture}
    
%     Probably the most important slide.
% \end{frame}

% \begin{frame}{Scope Freeze Milestones}
%     \begin{itemize}
%         \item A date where the scope cannot be altered
%         \item Typically, specifications are set
%         \item Also design freeze, customer approval freeze...
%         \item De-couple sponsor from product
%         \item De-couple project management from product?
%         \item There is agile!!
%     \end{itemize}
% \end{frame}

% \begin{frame}{MoScCoW - Scheme}
%     \begin{description}
%         \item[M] must have: non-negotiable
%         \item[S] should have: not vital, but adds significant value
%         \item[C] could have: nice to have but small impact when left out
%         \item[W] will not have: stuff without priority in the given time frame.
%     \end{description}    
% \end{frame}


\newcommand{\myax}[1]{
    \item 
    \begin{minipage}[t]{0.7\textwidth}
        #1
    \end{minipage}
}
\begin{frame}{Axioms of Communication (Watzlawick)}
    \begin{enumerate}
        \myax{ \textbf{One cannot not communicate;}\par\medskip every behaviour is a form of communication,\par\medskip  including body language and silence.}
        \end{enumerate}
\end{frame}

\begin{frame}{Axioms of Communication (Watzlawick)}\begin{enumerate}
        
        \myax{Every communication has a \textbf{content} and \textbf{relationship} aspect.\par\medskip The relationship aspect classifies the content aspect (metacommunication).\par\medskip Communication include more information than is present in the meanings of the words.}
        \end{enumerate}
\end{frame}

\begin{frame}{Axioms of Communication (Watzlawick)}\begin{enumerate}
        \myax{The punctuation (breaks, communication frequency) is key for the relationship aspect. \par\medskip
        Each person in a communication structure their communication flow differently and interprets their own behaviour as simply a reaction to the behaviour of the other person.\par\medskip
        Each event during a communication is seen as the cause of the next. (Most people are very unaware of their communication.)}
        \end{enumerate}
\end{frame}

\begin{frame}{Axioms of Communication (Watzlawick)}\begin{enumerate}
        
        \myax{Human communication uses multiple media and channels.\par\medskip (Communication is never a singular act.)}
        \end{enumerate}
\end{frame}

\begin{frame}{Axioms of Communication (Watzlawick)}\begin{enumerate}
        
        \myax{Communication can be \textbf{symmetric} or \textbf{not.} \par\medskip
        (Communication has to do with status and power.)
        }
    \end{enumerate}

    \vfill
    Source: \cite{watzlawick2000menschliche}
\end{frame}



\begin{frame}{Four-Sides-Model (Schulz von Thun)}
    \showmytikz{../shared/tikz/tikz-PM-four-sides-model}
    \vfill
    Source: \cite{von2013miteinander}
\end{frame}

\begin{frame}{Framing / "Sandwiching"}
    Positive - negative - positive
    
    \medskip
    
    \begin{quote}
        Dear xy,\\
        Thank you for... \\
        Unfortunately...\\
        We will do our best...\\
        Yours faithfully...
    \end{quote}
    
    \medskip

    \begin{quote}
        I like the approach, thank you for your work etc.\\
        I suggest next time / in additon we...\\
        Then we will reach perfection...\\
    \end{quote}

    Positive mindset and attitude! 

\end{frame}


\begin{frame}{Overcommunication}
    \begin{itemize}
        \item Communication should come to conclusions
        \item Decision making
        \item Proper meeting culture
        \item Talk until everything is said
        \item Talk more
    \end{itemize}
\end{frame}

\begin{frame}{Meeting culture}
    \begin{itemize}
        \item Have an agenda - in calendar/invitation
        \item Agree on agenda
        \item Follow agenda
        \item Project agenda (onto screen)
        \item Make a protocol with decisions
        \item[] - ideally by typing into agenda
        \item Make agenda for next meeting
        \item Distribute protocol - in calendar        
    \end{itemize}
\end{frame}
